\section{Introduction}

Semantic entropy, celebrated for capturing ``meaning-level uncertainty'' in large language models, underperforms a simple confidence score by 24 AUROC points on multiple-choice hallucination detection. This counterintuitive finding challenges the prevailing assumption that more sophisticated uncertainty quantification (UQ) methods necessarily yield better hallucination detection.

The ability to automatically detect when LLMs produce hallucinations---confident-sounding but factually incorrect outputs---remains a critical challenge for safe deployment. Uncertainty quantification offers a principled approach: outputs with high uncertainty may be unreliable and warrant additional verification. Recent advances have produced increasingly sophisticated UQ methods, from simple token entropy to semantic entropy with NLI-based clustering \cite{kuhn2023semantic}, self-consistency checks \cite{manakul2023selfcheckgpt}, and P(True) calibration \cite{kadavath2022language}.

However, a deeper problem emerges upon closer examination. Each method has been evaluated under favorable conditions using different benchmarks, data splits, and evaluation protocols. Kuhn et al.\ demonstrated semantic entropy's superiority over token entropy on their custom QA splits. Manakul et al.\ showed SelfCheckGPT's effectiveness on WikiBio generation. Kadavath et al.\ validated P(True) on proprietary data. Yet no controlled head-to-head comparison exists on identical conditions, leaving practitioners unable to select appropriate methods for their specific use cases.

We address this gap by conducting the first controlled comparison of UQ methods for hallucination detection on TruthfulQA's multiple-choice format. Our key insight is that \textbf{UQ method effectiveness depends critically on task format}: simple confidence-based methods extract discriminative signals directly from logits without format dependency, while semantic entropy's NLI-based clustering requires substantial semantic content that single-letter MC answers cannot provide. The semantic clustering mechanism works correctly---we verify that NLI models produce meaningful clusters with non-trivial entropy variance---but the resulting entropy scores do not correlate with hallucination status on MC-format tasks.

Building on this insight, we make the following contributions. First, we provide empirical evidence that token-level UQ methods (max probability, choice entropy) effectively discriminate hallucinations on TruthfulQA mc1, achieving AUROC 0.77--0.81 with Llama-3-8B. Second, we demonstrate that semantic entropy achieves only AUROC 0.56 on the same task despite mechanism verification confirming correct implementation. Third, we identify format dependency as the root cause: semantic entropy's design assumes free-form generation where NLI comparison is meaningful, while MC answers are single letters with no semantic content to compare. This finding has immediate practical implications: practitioners should match UQ method selection to task format rather than assuming more sophisticated methods perform better.

The remainder of this paper is organized as follows. Section~\ref{sec:related} positions our work relative to prior UQ methods. Section~\ref{sec:methodology} describes our methodology for controlled comparison. Section~\ref{sec:experiments} presents experimental setup, Section~\ref{sec:results} reports results, and Section~\ref{sec:discussion} discusses implications and limitations. Section~\ref{sec:conclusion} concludes with future directions.
